\subsection{A relative graph obstruction with the coefficient place free}

\begin{statement}{A.119}{Proved}
{Boundary splitting can obstruct even scalar-shifted relative repair}
Put
\[
 \begin{gathered}
 \alpha^9=3,\quad K=\Q(\alpha),\quad k=\Q(\alpha^3),\\
 T=\ker\bigl(\operatorname{Res}_{K/\Q}\mathbf G_m
            \xrightarrow{N_{K/k}}
                         \operatorname{Res}_{k/\Q}\mathbf G_m\bigr),
 \quad C=T[3],\\
 E_0=\Q(\mu_9,\alpha,(-2)^{1/3}),\qquad
                  q=321193,\quad g=5.
 \end{gathered}
 \tag{A.119.1}
\]
The prime \(q\) splits completely in \(E_0\), and \(g\)
is a primitive residue unit at \(q\).
Use the actual relative module
\[
 C=Z^3\F_3[Z]/(Z^9)\otimes\chi,\quad x=1+Z,\qquad
                         t=Z^3,\quad c=Z^4.
 \tag{A.119.2}
\]
Thus \(t=-e_0+e_3\) and
\(c=e_0-e_1-e_3+e_4\); both have zero sum on each
block of indices congruent modulo three.
Prescribe
\[
             \gamma_2=0,\qquad
       \gamma_q=t\,\kappa_q(q)+c\,\kappa_q(g).
 \tag{A.119.3}
\]
Leave \(\gamma_3\) unrestricted.

There is no global class \(\gamma\in H^1(\Q,C)\)
with (A.119.3) satisfying both of the following
requirements:
it is unramified outside \(\{2,3,q\}\) and finitely
many primes splitting completely in \(E_0\);
at every such new ramified prime \(p\), its tame
inertia and unit/Frobenius coordinates satisfy
\[
                  c_p=z_pt_p+w_p\mathbf1
                         \quad(z_p,w_p\in\F_3).
 \tag{A.119.4}
\]
Here \(\mathbf1=\sum_{i=0}^8e_i=Z^8\) is the
literal scalar coordinate direction.
In particular the assertion includes the cyclic
condition \(c_p=z_pt_p\).

The finite prescription (A.119.3) itself is realized
by actual global elements of \(T(\Q)\), with all
new Kummer ramification at primes splitting the
larger prior field
\[
 E=\Q\bigl(\mu_{216},3^{1/27},(-2)^{1/72},
       (\tau^i\beta_{\rm pre})^{1/3}:0\le i<9\bigr)
                       \supset E_0
 \tag{A.119.5}
\]
of A.106.  The same \(q\) splits this full field.
Thus the obstruction concerns the additional
equations (A.119.4), even with the place three
free and with literal lower norm one.
It does not assert that the specific discrepancy
of the ordinary ghost norm chain has the prescribed
local values (A.119.3).
\end{statement}

\begin{proof}
\emph{The rank-two form and its exact dual map.}
Use the nondegenerate form and its invariant operator
from A.113 or A.115:
\[
 \begin{gathered}
 B_*(v,w)=\sum_{h=0}^2(v_hw_{h+3}-v_{h+3}w_h),
 \qquad z=Z^2/(1+Z),\\
 F=z^2,\qquad B_0(v,w)=B_*(v,Fw).
 \end{gathered}
 \tag{A.119.6}
\]
All these are maps of the actual \(G_\Q\)-modules.
The form \(B_0:\bigwedge^2C\to\mu_3\) has rank two
and radical \(R=Z^5\F_3[Z]/Z^9\otimes\chi\).
Write
\[
 W=C/R,\qquad C'=Z^7\F_3[Z]/Z^9\otimes\chi.
 \tag{A.119.7}
\]
Multiplication by \(F\) induces an equivariant
isomorphism \(\bar F:W\longrightarrow C'\).
Indeed \(FZ^3=Z^7+Z^8\) and \(FZ^4=Z^8\).
In particular it sends the quotient socle \(Z^4\)
to the constant vector \(\mathbf1\).
The form \(B_0\) induces a nondegenerate alternating
form \(B_W\) on \(W\), and these displayed maps
satisfy the literal pairing identity
\[
 B_*(v,\bar Fw)=B_W(\bar v,w)
                         \quad(v\in C,\ w\in W).
 \tag{A.119.8}
\]
Also \(\mathbf1\in R\), so
\(B_0(v,\mathbf1)=0\) for every \(v\).
For the two vectors in (A.119.2),
\[
       B_0(t,c)=B_*(Z^3,Z^8)=(-1)-1=1
                                      \quad\hbox{in }\F_3.
 \tag{A.119.9}
\]

\emph{The prior boundary cocycle and the three-adic restriction.}
Let \(i:\mu_3\to C\) send one to \(\mathbf1\), and
put
\[
                  \delta=i_*\kappa_3(-2)\in H^1(\Q,C).
 \tag{A.119.10}
\]
This is an actual global cocycle factoring through
the splitting field \(\Q(\mu_3,(-2)^{1/3})\)
inside \(E_0\).
Its values lie in \(C'\).
Via \(\bar F\), it is the image of
\(\delta_W\in H^1(\Q,W)\), obtained by mapping
\(\kappa_3(-2)\) to the socle \(Z^4\) of \(W\).

We check that its localization at three is a
nonzero isotropic vector in a two-dimensional
cohomology space.
In the basis \(w_0=Z^3,w_1=Z^4\) of \(W\),
translation sends
\(w_0\mapsto w_0+w_1\), \(w_1\mapsto w_1\).
Complex conjugation fixes \(w_0\) and negates
\(w_1\).
The exact sequence is consequently
\[
 0\longrightarrow\mu_3
   \xrightarrow{\,1\mapsto w_1\,}W
   \xrightarrow{\,w_0\mapsto1\,}\F_3
   \longrightarrow0,
 \qquad \partial(1)=\kappa_3(3).
 \tag{A.119.11}
\]
The equality of the connecting class follows on
the actual radical generator and the involution
from the two displayed matrices; both cocycles
factor through the same \(S_3\) quotient.
The local group at three maps onto that quotient.

Now \(H^1(\Q_3,\mu_3)=\Q_3^\times/\Q_3^{\times3}\)
has dimension two.  Its valuation class
\(\kappa_3(3)\) and unit class
\(\kappa_3(-2)\) are independent:
the latter has zero valuation and is \(7\) modulo
nine, whereas unit cubes are \(1\) or \(-1\).
It follows from (A.119.11) that the image of
\(H^1(\Q_3,\mu_3)\) in \(H^1(\Q_3,W)\) is a
one-dimensional line generated by the nonzero
\(\delta_W\).
It is isotropic, since \(B_W\) vanishes on the
coefficient socle.

The module \(W\) and its Cartier dual have no
local invariants: their translation-fixed line
has involution sign \(-1\), and \(B_W\) identifies
them.
Local duality and the Euler characteristic
\cite[I, Corollary 2.3 and Theorem 2.8]{BIR_MilneADT}
therefore give
\[
 \dim_{\F_3}H^1(\Q_3,W)=2,\qquad
 \langle\ ,\ \rangle_{B_W}\ \hbox{perfect and symmetric}.
 \tag{A.119.12}
\]
The orthogonal of the nonzero isotropic line
\(\F_3\delta_W\) is that same line.
Its half-diagonal quadratic form is zero on it.

\emph{The first global pairing fixes the quadratic value at three.}
Suppose \(\gamma\) satisfies the claimed support
and old local conditions, temporarily without
using (A.119.4).
The cup product
\((B_*)_*(\gamma\smile\delta)\) is in
\(H^2(\Q,\mu_3)=\operatorname{Br}(\Q)[3]\).
At two its local invariant is zero by \(\gamma_2=0\).
At \(q\) and at every newly allowed prime,
\(\delta\) is locally zero, because those primes
split \(E_0\).
At any other finite prime except three both
classes are unramified, so their Brauer invariant
is zero.  The real three-primary group is zero.
Global reciprocity thus gives
\[
      \operatorname{inv}_3
          (B_*)_*(\gamma_3\smile\delta_3)=0.
 \tag{A.119.13}
\]
By (A.119.8), this is exactly
\(\langle\bar\gamma_3,\delta_W\rangle_{B_W}=0\).
Equations (A.119.11)--(A.119.12) imply
\(\bar\gamma_3\in\F_3\delta_W\), hence
\[
 \operatorname{inv}_3\left(
          \tfrac12(B_0)_*(\gamma_3\smile\gamma_3)\right)=0.
 \tag{A.119.14}
\]
This is the actual affine restriction imposed by
the prior boundary field; the value at three
was not prescribed in (A.119.3).

\emph{The second global pairing contradicts the new graphs.}
For any global class \(\gamma\) define
\[
       Q_0(\gamma)=\tfrac12(B_0)_*
                  (\gamma\smile\gamma)\in\operatorname{Br}(\Q)[3].
 \tag{A.119.15}
\]
At a split prime, a class with coordinates
\((t_p,c_p)\) has local invariant given, up to
the fixed nonzero uniformizer--unit factor, by
\(B_0(t_p,c_p)\).
This follows by expanding the cup product:
the two mixed terms agree, since both degree-one
cup product and \(B_0\) are alternating, and
the factor \(1/2\) cancels their sum.
For (A.119.4) this value is zero, because
\(B_0(t_p,t_p)=B_0(t_p,\mathbf1)=0\).
At any other unramified prime, including a
nonsplit one, the cup class inflates from the
unramified procyclic quotient, whose degree-two
cohomology with these finite torsion coefficients
is zero.  Its local invariant is therefore zero.
The local values at two and three vanish by
(A.119.3) and (A.119.14).

At \(q\), in contrast, (A.119.9) makes the value
that of \(\kappa_q(q)\smile\kappa_q(g)\).
It is nonzero: \(g\) is a primitive residue unit
and its cubic residue character is nontrivial.
Thus \(Q_0(\gamma)\) would have exactly one
nonzero local invariant, contradicting Brauer
reciprocity
\cite[I, Theorem 4.10(b)]{BIR_MilneADT}.
This proves the asserted nonexistence.

\emph{The explicit field and the compatible finite prescription.}
The data of A.106 give
\[
 \begin{gathered}
 q\equiv1\pmod{216},\quad
 2827^{27}\equiv3\pmod q,\quad
 (-2)^{(q-1)/72}\equiv1\pmod q,\\
 5^{(q-1)/3}\not\equiv1\pmod q.
 \end{gathered}
 \tag{A.119.16}
\]
Together with the certified conjugate characters
of its actual \(\beta_{\rm pre}\), these prove
that \(q\) splits (A.119.5), hence also \(E_0\).
The inclusions are literal:
raise \(3^{1/27}\) to its third power,
\((-2)^{1/72}\) to its twenty-fourth power,
and use \(\mu_9\subset\mu_{216}\).

At \(q\), the local relative torus is split.
The coordinates \(q^{t_i}g^{c_i}\) have product
one on every three-element block, so (A.119.3)
is the Kummer class of an actual relative
norm-one local point.
Apply A.117 with the full field (A.119.5),
including \(q\) in \(S\), prescribing (A.119.3)
and prescribing zero at every other member
of \(S\setminus\{3\}\).
It produces an actual global relative norm-one
element and permits new ramification only at
primes splitting (A.119.5).
Its zero local restrictions at the additional
old primes give no ramification there.
Therefore the finite plans are realized with
precisely the support allowed in the statement,
before (A.119.4) is imposed.
This proves the final assertion and its scope.
\end{proof}

The retained program
\path{code/norm_free_three_graph_certificate.py}
checks all the displayed coefficient maps, the
rank-two quotient extension, the relative witness,
the scalar-shifted graph identities, and the
explicit splitting and unit data.
